\documentclass[11pt,a4paper]{article}
% preamble.tex — shared setup for human-verification documents
% Compile each document with:  tectonic <File>.tex

\usepackage[T1]{fontenc}
\usepackage[utf8]{inputenc}
\usepackage{lmodern}
\usepackage{amssymb,amsmath,amsthm}
\usepackage{mathtools}
\usepackage[dvipsnames]{xcolor}
\usepackage{listings}
\usepackage{hyperref}
\usepackage[margin=2.5cm]{geometry}
\usepackage{parskip}

% ── Colours ──────────────────────────────────────────────────────────
\definecolor{leanBg}{HTML}{F7F7F0}
\definecolor{haskellBg}{HTML}{F0F4F7}
\definecolor{keyword}{HTML}{7B3294}
\definecolor{comment}{HTML}{5A8C32}
\definecolor{noteFg}{HTML}{B35900}

% ── Lean 4 listings language ─────────────────────────────────────────
% NOTE: Lean listing blocks use ASCII approximations of unicode symbols
% (e.g. -> for arrow, o. for compose).  The literate key substitutes
% these with proper math-mode renderings in the PDF.
\lstdefinelanguage{Lean4}{
  morekeywords={def,theorem,structure,where,namespace,end,import,fun,
                let,have,show,by,exact,intro,cases,rfl,subst,obtain,
                funext,simp,rw,apply,variable,Type,Prop,instance,
                section,open,class,extends,deriving,sorry,noncomputable},
  sensitive=true,
  morecomment=[l]{--},
  morecomment=[n]{/-}{-/},
  morestring=[b]",
}

% ── Lean listing style ───────────────────────────────────────────────
\lstdefinestyle{lean}{
  language=Lean4,
  backgroundcolor=\color{leanBg},
  basicstyle=\small\ttfamily,
  keywordstyle=\color{keyword}\bfseries,
  commentstyle=\color{comment}\itshape,
  breaklines=true,
  frame=single,
  framerule=0.4pt,
  xleftmargin=4pt,
  xrightmargin=4pt,
  aboveskip=6pt,
  belowskip=4pt,
  literate=
    {->}{{$\to$}}2
    {=>}{{$\Rightarrow$}}2
    {<|}{{$\langle$}}1
    {|>}{{$\rangle$}}1
    {==>}{{$\Longrightarrow$}}3
    { o. }{{$\;\circ\;$}}3
    {forall}{{$\forall$}}1
    {:=}{{$\,:=\,$}}2
    {!=}{{$\neq$}}2
    {<=}{{$\leq$}}2
    {>=}{{$\geq$}}2,
}

% ── Haskell listing style ───────────────────────────────────────────
\lstdefinestyle{haskell}{
  language=Haskell,
  backgroundcolor=\color{haskellBg},
  basicstyle=\small\ttfamily,
  keywordstyle=\color{keyword}\bfseries,
  commentstyle=\color{comment}\itshape,
  breaklines=true,
  frame=single,
  framerule=0.4pt,
  xleftmargin=4pt,
  xrightmargin=4pt,
  aboveskip=6pt,
  belowskip=4pt,
  literate=
    {->}{{$\to$}}2
    {=>}{{$\Rightarrow$}}2,
}

% ── Theorem environments ────────────────────────────────────────────
\theoremstyle{definition}
\newtheorem{defn}{Definition}[section]
\newtheorem{thm}[defn]{Theorem}
\newtheorem{prop}[defn]{Proposition}

\theoremstyle{remark}
\newtheorem*{fidelitynote}{Fidelity note}

% ── Convenience macros ───────────────────────────────────────────────
\newcommand{\lean}[1]{\texttt{#1}}      % inline Lean identifier
\newcommand{\hask}[1]{\texttt{#1}}      % inline Haskell identifier
\newcommand{\Ext}{\mathrm{Ext}}
\newcommand{\Pos}{\mathrm{Pos}}
\newcommand{\Shape}{\mathrm{Shape}}


\title{Human Verification: \texttt{Containers.DirichletClosed}\\[4pt]
  \large $(\mathbf{Cont}, \otimes, y)$ is a Closed Monoidal Category}
\author{Machine-generated from \texttt{lean/DirichletClosed.lean} for review by Neil Ghani}
\date{July 2026}

\begin{document}
\maketitle

\begin{abstract}
This document presents every definition and theorem \emph{statement} from the
Lean~4 module \lean{Containers.DirichletClosed} in three parallel forms: the
original Lean, a readable Haskell rendering, and precise mathematical English.
Proofs are omitted --- they are machine-checked by Lean.  The purpose is to let
a human verifier confirm that the formalisation captures the intended
mathematics.

The file proves that the Dirichlet tensor $\otimes$ on containers is
\textbf{closed}: for every container $q$ the functor $- \otimes q$ has a right
adjoint, the \emph{internal hom} $[q, r]$.  The adjunction is presented as an
explicit type equivalence $\lean{dirichlet\_closure}$; all round trips hold by
\lean{rfl}.  Naturality of currying (in both arguments) and the universal
property in counit form (evaluation and uniqueness) are also proved.  This is
Niu--Sp\'{e}ak Eq.~(4.79) in the container presentation.
\end{abstract}

\tableofcontents
\bigskip

%===================================================================
\section{Overview}
%===================================================================

The central objects are:
\begin{itemize}
  \item The \textbf{internal hom} $[q, r]$: a container whose shapes are the
    morphisms $q \Rightarrow r$ in $\mathbf{Cont}$.
  \item The \textbf{adjunction} $- \otimes q \dashv [q, -]$, given concretely
    as the currying bijection
    $(p \otimes q \Rightarrow r) \cong (p \Rightarrow [q, r])$.
\end{itemize}

The bijection is pure currying: the $q$-component of a backward position map
is packaged \emph{into the shape} of the internal hom (as its own backward
position map), and the $p$-component becomes the curried morphism's backward
position map.  This split is why the positions of $[q, r]$ are
$\Sigma\,t : S_q,\; r[u(t)]$ and not something else.

Every proof is \lean{rfl}: the operations involved are rearrangements of
$\Sigma/\times$ telescopes, and Lean 4's definitional $\eta$ for \lean{Prod},
\lean{Sigma}, and the \lean{ContainerMorphism} structure handles the rest.

%===================================================================
\section{Auxiliary: type equivalences}
%===================================================================

\subsection{Structure: \lean{TypeEquiv}}

\begin{defn}[Type equivalence]
A \textbf{type equivalence} $A \simeq B$ is a pair of maps with both round trips:
\[
  f : A \to B, \quad g : B \to A, \quad
  g \circ f = \mathrm{id}_A, \quad f \circ g = \mathrm{id}_B.
\]
This is a Mathlib-free stand-in for \lean{Equiv}.
\end{defn}

\paragraph{Lean.}
\begin{lstlisting}[style=lean]
structure TypeEquiv (A : Type u) (B : Type v) where
  toFun    : A -> B
  invFun   : B -> A
  left_inv : forall a, invFun (toFun a) = a
  right_inv: forall b, toFun (invFun b) = b
\end{lstlisting}

\paragraph{Haskell.}
\begin{lstlisting}[style=haskell]
data TypeEquiv a b = TypeEquiv
  { toFun    :: a -> b
  , invFun   :: b -> a
  , leftInv  :: forall x. invFun (toFun x) == x   -- propositional
  , rightInv :: forall y. toFun (invFun y) == y
  }
\end{lstlisting}


%===================================================================
\section{The internal hom}
%===================================================================

\subsection{Definition: \lean{Container.ihom}}

\begin{defn}[Internal hom for $\otimes$]
For containers $q$ and $r$, the \textbf{internal hom} $[q, r]$ is the container:
\begin{itemize}
  \item \emph{Shapes}: $\lean{ContainerMorphism}\,q\,r$ --- the hom-set of
    $\mathbf{Cont}$.  A shape is a pair $(u, \phi)$ with
    $u : S_q \to S_r$ and $\phi_t : r[u(t)] \to q[t]$ for each $t : S_q$.
  \item \emph{Positions} over shape $f = (u, \phi)$:
    $\Sigma\,t : S_q,\; r[u(t)]$.
\end{itemize}
\end{defn}

\paragraph{Lean.}
\begin{lstlisting}[style=lean]
def Container.ihom (q r : Container) : Container where
  Shape := ContainerMorphism q r
  Pos   := fun f => (t : q.Shape) * r.Pos (f.onShapes t)
\end{lstlisting}

\paragraph{Haskell.}
\begin{lstlisting}[style=haskell]
-- Internal hom [q, r]:
-- Shape  = ContainerMorphism q r
-- Pos f  = Exists t. r.Pos (onShapes f t)   (dependent pair)
data IhomPos q r f = IhomPos { iqShape :: Shape q, irPos :: Pos r (onShapes f iqShape) }
-- (Haskell cannot express the dependent type directly.)
\end{lstlisting}

\begin{fidelitynote}
The fact that the \emph{shape} type of $[q, r]$ is literally \lean{ContainerMorphism
q r} is not a notational convenience --- it is the mathematical content
(Niu--Sp\'{e}ak Eq.~4.79).  The positions over $f$ are pairs $(t, i)$ where
$t : S_q$ and $i : r[f(t)]$ is an $r$-position over the image shape.  In
Haskell, the dependent type $r.\mathrm{Pos}(f.\lean{onShapes}(t))$ cannot be
expressed directly; the Haskell rendering is therefore approximate.
\end{fidelitynote}


\subsection{Theorem: \lean{Container.ihom\_shape\_eq}}

\begin{thm}[Shape of the internal hom]
The shape type of $[q, r]$ is definitionally the hom-set:
\[
  [q, r].\mathrm{Shape} = \mathrm{ContainerMorphism}\,q\,r.
\]
\end{thm}

\paragraph{Lean.}
\begin{lstlisting}[style=lean]
theorem Container.ihom_shape_eq (q r : Container) :
    (Container.ihom q r).Shape = ContainerMorphism q r := rfl
\end{lstlisting}


\subsection{Theorem: \lean{Container.ihom\_pos\_eq}}

\begin{thm}[Positions of the internal hom]
The position type of $[q, r]$ over a shape $f : q \Rightarrow r$ is:
\[
  [q, r].\mathrm{Pos}\,f = \Sigma\,t : S_q,\; r[f(t)].
\]
\end{thm}

\paragraph{Lean.}
\begin{lstlisting}[style=lean]
theorem Container.ihom_pos_eq (q r : Container)
    (f : ContainerMorphism q r) :
    (Container.ihom q r).Pos f
      = ((t : q.Shape) * r.Pos (f.onShapes t)) := rfl
\end{lstlisting}


%===================================================================
\section{Currying and uncurrying}
%===================================================================

\subsection{Definition: \lean{dirCurry}}

\begin{defn}[Currying]
Given $F : p \otimes q \Rightarrow r$, define
$\lean{dirCurry}(F) : p \Rightarrow [q, r]$ by:
\begin{itemize}
  \item \emph{On shapes}: $s \mapsto \{$shape map: $t \mapsto F(s,t)$;
    backward position map: $t, i \mapsto (\lean{F.onPos}(s,t,i)).2\}$.
    (The $q$-component of $F$'s backward position map becomes the shape's own
    backward position map.)
  \item \emph{On positions} (backward): $(s, (t, i)) \mapsto
    (\lean{F.onPos}(s,t,i)).1$.
    (The $p$-component of $F$'s backward position map.)
\end{itemize}
\end{defn}

\paragraph{Lean.}
\begin{lstlisting}[style=lean]
def dirCurry {p q r : Container}
    (F : ContainerMorphism (p |*| q) r) :
    ContainerMorphism p (Container.ihom q r) where
  onShapes := fun s =>
    { onShapes := fun t   => F.onShapes (s, t)
      onPos    := fun t i => (F.onPos (s, t) i).2 }
  onPos := fun s x => (F.onPos (s, x.1) x.2).1
\end{lstlisting}

\paragraph{Haskell.}
\begin{lstlisting}[style=haskell]
dirCurry :: ContainerMorphism (p `Dir` q) r
         -> ContainerMorphism p (Ihom q r)
dirCurry f = Morph
  { onShapes = \s -> Morph
      { onShapes = \t   -> onShapes f (s, t)
      , onPos    = \t i -> snd (onPos f (s, t) i) }
  , onPos    = \s (t, i) -> fst (onPos f (s, t) i) }
\end{lstlisting}


\subsection{Definition: \lean{dirUncurry}}

\begin{defn}[Uncurrying]
Given $G : p \Rightarrow [q, r]$, define
$\lean{dirUncurry}(G) : p \otimes q \Rightarrow r$ by:
\begin{itemize}
  \item \emph{On shapes}: $(s, t) \mapsto (G(s)).(t)$, i.e.\ apply the
    morphism $G(s) : q \Rightarrow r$ to the shape $t$.
  \item \emph{On positions} (backward): $(s,t), i \mapsto
    (G.\lean{onPos}(s, (t, i)),\;(G(s)).\lean{onPos}(t, i))$,
    producing a pair of a $p$-position and a $q$-position.
\end{itemize}
\end{defn}

\paragraph{Lean.}
\begin{lstlisting}[style=lean]
def dirUncurry {p q r : Container}
    (G : ContainerMorphism p (Container.ihom q r)) :
    ContainerMorphism (p |*| q) r where
  onShapes := fun st => (G.onShapes st.1).onShapes st.2
  onPos    := fun st i =>
    (G.onPos st.1 <|st.2, i|>, (G.onShapes st.1).onPos st.2 i)
\end{lstlisting}

\paragraph{Haskell.}
\begin{lstlisting}[style=haskell]
dirUncurry :: ContainerMorphism p (Ihom q r)
           -> ContainerMorphism (p `Dir` q) r
dirUncurry g = Morph
  { onShapes = \(s, t) -> onShapes (onShapes g s) t
  , onPos    = \(s, t) i ->
      (onPos g s (t, i), onPos (onShapes g s) t i) }
\end{lstlisting}


%===================================================================
\section{Round trips and the main theorem}
%===================================================================

\subsection{Theorem: \lean{dirUncurry\_dirCurry}}

\begin{thm}[Uncurrying undoes currying]
For all $F : p \otimes q \Rightarrow r$:
\[
  \lean{dirUncurry}(\lean{dirCurry}(F)) = F.
\]
Holds by \lean{rfl}.
\end{thm}

\paragraph{Lean.}
\begin{lstlisting}[style=lean]
theorem dirUncurry_dirCurry {p q r : Container}
    (F : ContainerMorphism (p |*| q) r) :
    dirUncurry (dirCurry F) = F := rfl
\end{lstlisting}

\paragraph{Haskell.}
\begin{lstlisting}[style=haskell]
-- dirUncurry (dirCurry f) = f   -- by eta and product eta
\end{lstlisting}


\subsection{Theorem: \lean{dirCurry\_dirUncurry}}

\begin{thm}[Currying undoes uncurrying]
For all $G : p \Rightarrow [q, r]$:
\[
  \lean{dirCurry}(\lean{dirUncurry}(G)) = G.
\]
Holds by \lean{rfl}.
\end{thm}

\paragraph{Lean.}
\begin{lstlisting}[style=lean]
theorem dirCurry_dirUncurry {p q r : Container}
    (G : ContainerMorphism p (Container.ihom q r)) :
    dirCurry (dirUncurry G) = G := rfl
\end{lstlisting}

\paragraph{Haskell.}
\begin{lstlisting}[style=haskell]
-- dirCurry (dirUncurry g) = g   -- by eta for ContainerMorphism structure
\end{lstlisting}


\subsection{Definition: \lean{dirichlet\_closure}}

\begin{defn}[The Dirichlet tensor is closed]
For all containers $p, q, r$, there is a type equivalence:
\[
  (p \otimes q \Rightarrow r) \;\simeq\; (p \Rightarrow [q, r]).
\]
This is the adjunction $- \otimes q \dashv [q, -]$, or equivalently,
$(\mathbf{Cont}, \otimes, y)$ is a \textbf{closed monoidal category}.

This is Niu--Sp\'{e}ak, \emph{Polynomial Functors}, Eq.~(4.79), in the
container presentation.

Declared as a \lean{def} (not a \lean{theorem}) because a type equivalence is
data (two maps plus two round-trip proofs).
\end{defn}

\paragraph{Lean.}
\begin{lstlisting}[style=lean]
def dirichlet_closure (p q r : Container) :
    ContainerMorphism (p |*| q) r <=> ContainerMorphism p (Container.ihom q r) where
  toFun    := dirCurry
  invFun   := dirUncurry
  left_inv := dirUncurry_dirCurry
  right_inv:= dirCurry_dirUncurry
\end{lstlisting}

\paragraph{Haskell.}
\begin{lstlisting}[style=haskell]
dirichletClosure :: TypeEquiv
  (ContainerMorphism (p `Dir` q) r)
  (ContainerMorphism p (Ihom q r))
dirichletClosure = TypeEquiv
  { toFun    = dirCurry
  , invFun   = dirUncurry
  , leftInv  = dirUncurryCurry
  , rightInv = dirCurryUncurry }
\end{lstlisting}


%===================================================================
\section{Naturality}
%===================================================================

\subsection{Definition: \lean{ihomMap}}

\begin{defn}[Functorial action of the internal hom]
For $\psi : r \Rightarrow r'$, the action
$\lean{ihomMap}(\psi) : [q, r] \Rightarrow [q, r']$:
\begin{itemize}
  \item On shapes: postcompose, $f \mapsto f \cdot \psi$.
    (A shape of $[q,r]$ is a morphism $q \Rightarrow r$; postcompose with $\psi$
    to get a morphism $q \Rightarrow r'$.)
  \item On positions (backward): given a position $(t, i)$ of $[q, r']$ over
    $f \cdot \psi$ (where $i : r'[\psi(f(t))]$), push $i$ back through
    $\psi^\sharp$ to get an $r$-position $\psi^\sharp_{f(t)}(i)$, giving
    $(t,\, \psi^\sharp_{f(t)}(i))$ of $[q, r]$.
\end{itemize}
\end{defn}

\paragraph{Lean.}
\begin{lstlisting}[style=lean]
def ihomMap {q r r' : Container} (psi : ContainerMorphism r r') :
    ContainerMorphism (Container.ihom q r) (Container.ihom q r') where
  onShapes := fun f => f.comp psi
  onPos    := fun f x => <|x.1, psi.onPos (f.onShapes x.1) x.2|>
\end{lstlisting}

\paragraph{Haskell.}
\begin{lstlisting}[style=haskell]
ihomMap :: ContainerMorphism r r'
        -> ContainerMorphism (Ihom q r) (Ihom q r')
ihomMap psi = Morph
  { onShapes = \f -> f `compM` psi
  , onPos    = \f (t, i) -> (t, onPos psi (onShapes f t) i) }
\end{lstlisting}


\subsection{Theorem: \lean{ihomMap\_id}}

\begin{thm}[Internal hom preserves identities]
$\lean{ihomMap}(\mathrm{id}_r) = \mathrm{id}_{[q,r]}$.  Holds by \lean{rfl}.
\end{thm}

\paragraph{Lean.}
\begin{lstlisting}[style=lean]
theorem ihomMap_id (q r : Container) :
    ihomMap (q := q) (ContainerMorphism.id r)
      = ContainerMorphism.id (Container.ihom q r) := rfl
\end{lstlisting}


\subsection{Theorem: \lean{ihomMap\_comp}}

\begin{thm}[Internal hom preserves composition]
$\lean{ihomMap}(\psi \cdot \psi') = \lean{ihomMap}(\psi) \cdot \lean{ihomMap}(\psi')$.
Holds by \lean{rfl}.
\end{thm}

\paragraph{Lean.}
\begin{lstlisting}[style=lean]
theorem ihomMap_comp {q r r' r'' : Container}
    (psi : ContainerMorphism r r') (psi' : ContainerMorphism r' r'') :
    ihomMap (q := q) (psi.comp psi')
      = (ihomMap psi).comp (ihomMap psi') := rfl
\end{lstlisting}


\subsection{Theorem: \lean{dirCurry\_naturality\_left}}

\begin{thm}[Naturality in $p$ (currying direction)]
For $h : p' \Rightarrow p$ and $F : p \otimes q \Rightarrow r$:
\[
  \lean{dirCurry}((h \otimes \mathrm{id}_q) \cdot F)
  \;=\;
  h \cdot \lean{dirCurry}(F).
\]
Holds by \lean{rfl}.
\end{thm}

\paragraph{Lean.}
\begin{lstlisting}[style=lean]
theorem dirCurry_naturality_left {p p' q r : Container}
    (h : ContainerMorphism p' p)
    (F : ContainerMorphism (p |*| q) r) :
    dirCurry ((Container.dir2 h (ContainerMorphism.id q)).comp F)
      = h.comp (dirCurry F) := rfl
\end{lstlisting}

\paragraph{Haskell.}
\begin{lstlisting}[style=haskell]
-- naturality_left: dirCurry ((h `tensor` id) `compM` f) = h `compM` dirCurry f
\end{lstlisting}


\subsection{Theorem: \lean{dirCurry\_naturality\_right}}

\begin{thm}[Naturality in $r$ (currying direction)]
For $F : p \otimes q \Rightarrow r$ and $\psi : r \Rightarrow r'$:
\[
  \lean{dirCurry}(F \cdot \psi)
  \;=\;
  \lean{dirCurry}(F) \cdot [q, \psi].
\]
Holds by \lean{rfl}.
\end{thm}

\paragraph{Lean.}
\begin{lstlisting}[style=lean]
theorem dirCurry_naturality_right {p q r r' : Container}
    (F : ContainerMorphism (p |*| q) r)
    (psi : ContainerMorphism r r') :
    dirCurry (F.comp psi) = (dirCurry F).comp (ihomMap psi) := rfl
\end{lstlisting}

\paragraph{Haskell.}
\begin{lstlisting}[style=haskell]
-- naturality_right: dirCurry (f `compM` psi) = dirCurry f `compM` ihomMap psi
\end{lstlisting}


\subsection{Theorem: \lean{dirUncurry\_naturality\_left}}

\begin{thm}[Naturality in $p$ (uncurrying direction)]
For $h : p' \Rightarrow p$ and $G : p \Rightarrow [q, r]$:
\[
  \lean{dirUncurry}(h \cdot G)
  \;=\;
  (h \otimes \mathrm{id}_q) \cdot \lean{dirUncurry}(G).
\]
Holds by \lean{rfl}.
\end{thm}

\paragraph{Lean.}
\begin{lstlisting}[style=lean]
theorem dirUncurry_naturality_left {p p' q r : Container}
    (h : ContainerMorphism p' p)
    (G : ContainerMorphism p (Container.ihom q r)) :
    dirUncurry (h.comp G)
      = (Container.dir2 h (ContainerMorphism.id q)).comp
          (dirUncurry G) := rfl
\end{lstlisting}


%===================================================================
\section{The counit: evaluation}
%===================================================================

\subsection{Definition: \lean{dirEval}}

\begin{defn}[Evaluation morphism]
The \textbf{counit} of the adjunction $- \otimes q \dashv [q,-]$ is the
\textbf{evaluation} morphism
\[
  \mathrm{ev} : [q, r] \otimes q \Rightarrow r,
\]
obtained by uncurrying the identity on $[q, r]$.  Explicitly:
\begin{itemize}
  \item On shapes: $(f, t) \mapsto f(t)$, i.e.\ \emph{apply} the morphism $f$
    stored in the first component to the shape $t$.
  \item On positions (backward): $i \mapsto ((t, i),\; f^\sharp_t(i))$,
    splitting an $r$-position $i$ into a $[q,r]$-position and a $q$-position.
\end{itemize}
\end{defn}

\paragraph{Lean.}
\begin{lstlisting}[style=lean]
def dirEval (q r : Container) :
    ContainerMorphism (Container.ihom q r |*| q) r :=
  dirUncurry (ContainerMorphism.id (Container.ihom q r))
\end{lstlisting}

\paragraph{Haskell.}
\begin{lstlisting}[style=haskell]
dirEval :: ContainerMorphism (Ihom q r `Dir` q) r
dirEval = dirUncurry idMorph
-- Concretely:
--   onShapes (f, t) = onShapes f t
--   onPos (f, t) i  = ((t, i), onPos f t i)
\end{lstlisting}


\subsection{Theorem: \lean{dirEval\_dirCurry}}

\begin{thm}[Universal property: existence]
For every $F : p \otimes q \Rightarrow r$:
\[
  (\lean{dirCurry}(F) \otimes \mathrm{id}_q) \cdot \mathrm{ev} \;=\; F.
\]
Holds by \lean{rfl}.  This is the existence half of the universal property of
the counit.
\end{thm}

\paragraph{Lean.}
\begin{lstlisting}[style=lean]
theorem dirEval_dirCurry {p q r : Container}
    (F : ContainerMorphism (p |*| q) r) :
    (Container.dir2 (dirCurry F) (ContainerMorphism.id q)).comp
        (dirEval q r) = F := rfl
\end{lstlisting}

\paragraph{Haskell.}
\begin{lstlisting}[style=haskell]
-- (dirCurry f `tensor` id) `compM` dirEval = f
\end{lstlisting}


\subsection{Theorem: \lean{dirCurry\_unique}}

\begin{thm}[Universal property: uniqueness]
If $G : p \Rightarrow [q, r]$ satisfies
$(G \otimes \mathrm{id}_q) \cdot \mathrm{ev} = F$,
then $G = \lean{dirCurry}(F)$.
\end{thm}

\paragraph{Lean.}
\begin{lstlisting}[style=lean]
theorem dirCurry_unique {p q r : Container}
    (F : ContainerMorphism (p |*| q) r)
    (G : ContainerMorphism p (Container.ihom q r))
    (h : (Container.dir2 G (ContainerMorphism.id q)).comp
           (dirEval q r) = F) :
    G = dirCurry F
\end{lstlisting}

\paragraph{Haskell.}
\begin{lstlisting}[style=haskell]
-- uniqueness: if (g `tensor` id) `compM` ev = f, then g = dirCurry f.
\end{lstlisting}

\begin{fidelitynote}
The uniqueness proof is the only non-\lean{rfl} proof in the file: it uses
\lean{subst} to substitute the hypothesis $h$, reducing to
$\lean{dirCurry}(\lean{dirUncurry}(G)) = G$, which is exactly
\lean{dirCurry\_dirUncurry} read backwards.  So the mathematical content is
still the round-trip, applied once.
\end{fidelitynote}

\end{document}
